\documentclass[11pt,a4paper]{article}

% ---------------------------------------------------------
% PREAMBLE (stable, warning-free, Overleaf-safe)
% ---------------------------------------------------------
\usepackage[utf8]{inputenc}
\usepackage[T1]{fontenc}
\usepackage{lmodern}
\usepackage{amsmath, amssymb, amsfonts}
\usepackage{bm}
\usepackage{geometry}
\usepackage{graphicx}
\usepackage{hyperref}
\usepackage{cite}
\usepackage{authblk}
\usepackage{physics}
\usepackage{mathtools}

\geometry{margin=2.4cm}

\title{\textbf{Companion XXXV: The Galaxy–Halo Connection in the Relaxation-Driven Cyclic Cosmology}}
\author{Michael Lehmann \\ \texttt{mi.lehmann@gmx.de}}
\date{April 2026}

\begin{document}
\maketitle

\begin{abstract}
Galaxies form within dark matter halos, yet the relation between the luminous and 
the dark components is neither trivial nor universal. In the standard cosmological 
model, the galaxy–halo connection is governed by the hierarchical growth of halos, 
the efficiency of gas cooling, and the feedback processes that regulate star 
formation. In the Relaxation-Driven Cyclic Cosmology (RDCC), the scale-dependent 
evolution of the gravitational potential alters the timing and efficiency of halo 
collapse, thereby reshaping the environments in which galaxies form. This Companion 
develops a continuous, dynamical description of the galaxy–halo connection in RDCC, 
deriving how the relaxation scale influences stellar mass assembly, star formation 
efficiency, and the mapping between halo mass and galaxy luminosity. The resulting 
framework provides a natural explanation for several observed anomalies in galaxy 
statistics and offers new predictions for upcoming surveys.
\end{abstract}
% ---------------------------------------------------------
\section{Introduction}

The connection between galaxies and their host halos is one of the central pillars 
of modern cosmology. It links the invisible scaffolding of dark matter to the 
observable Universe of stars, gas, and radiation. In the standard cosmological 
model, this connection is often described through empirical relations such as the 
stellar-to-halo mass relation or the halo occupation distribution. These relations 
encode the cumulative effects of gas cooling, star formation, feedback, and 
environmental processes.

In RDCC, the situation is more intricate. The relaxation-modified potential 
evolution introduces a scale-dependent temporal structure into the growth of halos. 
This affects not only when halos form but also the dynamical environment in which 
gas accretes and cools. As a result, the galaxy–halo connection becomes sensitive 
to the relaxation scale $k_{\mathrm{rel}}$, which imprints itself on the efficiency 
of star formation and the luminosity of galaxies across cosmic time.
% ---------------------------------------------------------
\section{Halo Growth and Gas Accretion in RDCC}

The growth of a halo is governed by the evolution of its gravitational potential. 
In RDCC, this potential takes the form
\begin{equation}
    \Phi_{\mathrm{RDCC}}(k,a)
    = \Phi(k,a)
      \left[1 - \chi_{\Phi}
      \left(\frac{k}{k_{\mathrm{rel}}}\right)^{\alpha}\right].
\end{equation}

The suppression of small-scale modes delays the collapse of low-mass halos, while 
the enhanced coherence of large-scale modes accelerates the growth of massive halos. 
This asymmetry modifies the gas accretion rate. For a halo of mass $M$, the 
accretion rate can be approximated as
\begin{equation}
    \dot{M}_{\mathrm{gas}}(M,a)
    \propto \frac{\partial}{\partial a}
    \left[
        M(a)
        \left(
            1 - \chi_{\mathrm{acc}}
            \left(\frac{k(M)}{k_{\mathrm{rel}}}\right)^{\alpha}
        \right)
    \right].
\end{equation}

Halos near the relaxation scale experience the strongest modulation, as their 
growth is most sensitive to the interplay between suppressed small-scale modes 
and coherent large-scale modes.
% ---------------------------------------------------------
\section{Star Formation Efficiency and the Relaxation Scale}

Star formation is regulated by the competition between gas cooling and feedback. 
In RDCC, the altered collapse history of halos changes the density and temperature 
structure of the infalling gas. The star formation efficiency $\epsilon_{\star}$ 
can be expressed as
\begin{equation}
    \epsilon_{\star}(M,a)
    = \epsilon_{0}(M,a)
      \left[
        1 - \chi_{\star}
        \left(\frac{k(M)}{k_{\mathrm{rel}}}\right)^{\alpha}
      \right],
\end{equation}
where $\epsilon_{0}$ is the efficiency in the standard cosmological model.

Low-mass halos, whose collapse is delayed by the relaxation mechanism, exhibit 
reduced star formation efficiency. Massive halos, which form earlier and in more 
coherent environments, maintain or even enhance their efficiency. This produces a 
stellar-to-halo mass relation that bends more sharply at low masses and flattens 
at high masses, in agreement with several observational trends.
% ---------------------------------------------------------
\section{The Stellar-to-Halo Mass Relation in RDCC}

The stellar mass $M_{\star}$ of a galaxy is the time integral of its star formation 
rate. In RDCC, this becomes
\begin{equation}
    M_{\star}(M)
    = \int da\,
      \dot{M}_{\mathrm{gas}}(M,a)\,
      \epsilon_{\star}(M,a).
\end{equation}

Substituting the RDCC-modified expressions yields
\begin{equation}
    M_{\star,\mathrm{RDCC}}(M)
    = M_{\star}(M)
      \left[
        1 - \chi_{\mathrm{SHMR}}
        \left(\frac{k(M)}{k_{\mathrm{rel}}}\right)^{\alpha}
      \right].
\end{equation}

This relation naturally produces a suppressed stellar mass for dwarf-scale halos 
and a nearly unchanged stellar mass for cluster-scale halos. The transition mass 
corresponds to the relaxation scale and provides a direct observational probe of 
RDCC.
% ---------------------------------------------------------
\section{Galaxy Luminosity and the RDCC Potential}

The luminosity of a galaxy depends on its stellar mass and star formation history. 
In RDCC, the altered collapse sequence and gas accretion history lead to a 
luminosity function that deviates from the standard Schechter form. The faint-end 
slope becomes shallower due to the suppressed formation of low-mass halos, while 
the bright end remains largely unchanged.

The luminosity $L$ can be approximated as
\begin{equation}
    L_{\mathrm{RDCC}}(M)
    \propto M_{\star,\mathrm{RDCC}}(M)^{\gamma},
\end{equation}
where $\gamma$ encodes the stellar population properties.

This framework provides a natural explanation for the observed scarcity of 
low-luminosity galaxies and the early appearance of massive galaxies in deep 
surveys.
% ---------------------------------------------------------
\section{Conclusion}

The galaxy–halo connection in the Relaxation-Driven Cyclic Cosmology reflects the 
scale-dependent evolution of the gravitational potential. The relaxation mechanism 
modifies halo growth, gas accretion, star formation efficiency, and the mapping 
between halo mass and galaxy luminosity. These effects produce a distinctive 
stellar-to-halo mass relation, a modified luminosity function, and a natural 
explanation for several observed anomalies in galaxy statistics. The present 
Companion establishes the foundation for future studies of galaxy evolution, 
feedback processes, and the interplay between baryons and dark matter in RDCC.

% ========================
% References
% ========================

\begin{thebibliography}{10}

\bibitem{Flagship} Lehmann, M. (2026). Relaxation-Driven Cyclic Cosmology (RDCC): A Minimal CPT-Symmetric Framework with a Single Parameter. Flagship Paper. Zenodo: \url{https://zenodo.org/records/19744958} (v43.0 or later).

\bibitem{Zenodo} The complete RDCC companion series (I--LVII) is available at the same Zenodo record above.

% Thematisch relevante Companions
\bibitem{CompXVI} Companion XVI: Boltzmann Code and Modified Hierarchy.
\bibitem{CompXXXI} Companion XXXI: RDCC Halo Model and Non-Linear Structure.
\bibitem{CompXXXII} Companion XXXII--XXXIX: Galaxy--Halo Connection, Cosmic Web, Lensing, RSD, ISW, tSZ, Rotation Curves, CGM.

% Wichtige externe Referenzen
\bibitem{Planck2020} Planck Collaboration (2020), Astron. Astrophys. 641, A6.
\bibitem{DESI2024} DESI Collaboration (2024/2025), arXiv: [aktuelle $f\sigma_8$/BAO-Paper].
\bibitem{JWST2024} Maiolino et al. (2024), Nature (JWST high-$z$ galaxies/quasars).
\bibitem{Kaiser1987} Kaiser, N. (1987), Mon. Not. R. Astron. Soc. 227, 1 (RSD).
\bibitem{Bartelmann2001} Bartelmann, M. \& Schneider, P. (2001), Phys. Rep. 340, 291 (Weak Lensing Review).
\bibitem{HuOkamoto2002} Hu, W. \& Okamoto, T. (2002), Astrophys. J. 574, 566 (CMB Lensing).
\bibitem{LewisChallinor2006} Lewis, A. \& Challinor, A. (2006), Phys. Rep. 429, 1 (CMB Lensing Review).
\bibitem{NFW1997} Navarro, J. F., Frenk, C. S. \& White, S. D. M. (1997), Astrophys. J. 490, 493 (Halo Profiles).

\end{thebibliography}

\end{document}
